\section{Reward Formulation}
\label{sec:rewards}
\setcounter{table}{0}
\renewcommand{\thetable}{A.\arabic{table}}

\Cref{tab:rew_params} lists the weights of the total reward of
\cref{eq:total-reward} together with the hyperparameters of the individual
terms.

\begin{table}[H]
    \centering
    \caption{Reward weights and hyperparameters}
    \label{tab:rew_params}
    \setlength{\tabcolsep}{4pt}
    \begin{tabular}{@{}llr@{}}
        \toprule
        \textbf{Description} & \textbf{Symbol} & \textbf{Value} \\
        \midrule
        \multicolumn{3}{@{}l}{\textit{Term weights}} \\
        Pose tracking          & $\omega_{pose}$      & 5.0 \\
        Stability              & $\omega_{s}$         & 0.4 \\
        Leg power              & $\omega_{legs}$      & 0.5 \\
        Arm power              & $\omega_{arm}$       & 0.2 \\
        Leg action smoothness  & $\omega_{act,legs}$  & 0.2 \\
        Arm action smoothness  & $\omega_{act,arm}$   & 0.25 \\
        \addlinespace
        \multicolumn{3}{@{}l}{\textit{Term hyperparameters}} \\
        Base proximity threshold     & $\beta$          & 0.3 m \\
        Base proximity sensitivity   & $\sigma_{base}$  & 0.5 \\
        Leg power normalisation      & $\mu_{legs}$     & 450 W \\
        Arm power normalisation      & $\mu_{arm}$      & 20 W \\
        Action smoothness sensitivity& $\sigma_{act}$   & 0.25 \\
        \bottomrule
    \end{tabular}
\end{table}

The kernel widths of \cref{eq:pos-rot-rew} are piecewise-constant functions
of the instantaneous tracking error, selected at each step by the interval containing
that error, as listed in \cref{tab:sigma_schedule}.

\begin{table}[H]
    \centering
    \caption{Error-adaptive kernel width schedules}
    \label{tab:sigma_schedule}
    \setlength{\tabcolsep}{4pt}
    \begin{tabular}{@{}cc@{\hskip 1.5em}cc@{}}
        \toprule
        \multicolumn{2}{@{}c}{\textbf{Position}} & \multicolumn{2}{c@{}}{\textbf{Orientation}} \\
        \cmidrule(r){1-2}\cmidrule(l){3-4}
        $e_{\text{pos}}$ [m] & $\sigma_{\text{pos}}$ & $e_{\text{rot}}$ [rad] & $\sigma_{\text{rot}}$ \\
        \midrule
        $[0.75,\infty)$ & 0.500 & $[1.00,\infty)$ & 0.50 \\
        $[0.50, 0.75)$  & 0.400 & $[0.50, 1.00)$  & 0.40 \\
        $[0.25, 0.50)$  & 0.250 & $[0.25, 0.50)$  & 0.25 \\
        $[0.10, 0.25)$  & 0.150 & $[0.10, 0.25)$  & 0.15 \\
        $[0.05, 0.10)$  & 0.050 & $[0.00, 0.10)$  & 0.05 \\
        $[0.00, 0.05)$  & 0.025 & ---             & ---  \\
        \bottomrule
    \end{tabular}
\end{table}

\section{Task and Environment Configuration}
\label{sec:terrain}
\setcounter{table}{0}
\renewcommand{\thetable}{B.\arabic{table}}

\Cref{tab:env_config} reports the sampling ranges of the commanded end-effector
pose, expressed in the base frame, and the parameters governing procedural terrain
generation at each difficulty level. Command positions are expanded linearly from the
initial to the final range over the course of the command curriculum, while
orientations are sampled from the stated ranges throughout training.

\begin{table}[H]
    \centering
    \caption{Command sampling ranges and terrain generation parameters}
    \label{tab:env_config}
    \setlength{\tabcolsep}{4pt}
    \begin{tabular}{@{}lcc@{}}
        \toprule
        \multicolumn{3}{@{}l}{\textit{Command ranges}} \\
        \textbf{Component} & \textbf{Initial} & \textbf{Final} \\
        \midrule
        $x$ [m] & $[-0.30,\ 0.30]$ & $[-0.30,\ 2.50]$ \\
        $y$ [m] & $[-0.30,\ 0.30]$ & $[-1.00,\ 1.00]$ \\
        $z$ [m] & $[-0.10,\ 0.10]$ & $[-0.35,\ 0.20]$ \\
        Roll [rad]  & \multicolumn{2}{c@{}}{$0$} \\
        Pitch [rad] & \multicolumn{2}{c@{}}{$[-\pi/4,\ \pi/2]$} \\
        Yaw [rad]   & \multicolumn{2}{c@{}}{$[-\pi/2,\ \pi/2]$} \\
        \addlinespace
        \multicolumn{3}{@{}l}{\textit{Terrain parameters}} \\
        \textbf{Terrain} & \textbf{Level 1} & \textbf{Level 5} \\
        \midrule
        Rough, noise scale [m]  & 0.01 & 0.16 \\
        Hills, amplitude [m]    & 0.10 & 0.50 \\
        Craters, slope [deg]    & 10   & 30 \\
        \bottomrule
    \end{tabular}
\end{table}

Intermediate difficulty levels are obtained by linear interpolation between the
extremes reported above.

\section{Constraint Limits}
\setcounter{table}{0}
\renewcommand{\thetable}{C.\arabic{table}}

\Cref{tab:constraint_lims} reports the limits $l_i$ of \cref{eq:generic-constraint}
for each constraint. When a monitored quantity exceeds its limit, the CaT mechanism
induces a termination of the reward stream with a probability proportional to the
magnitude of the violation and capped at $p_i^{\max}$.

\begin{table}[H]
    \centering
    \caption{Constraint limits. Soft constraints use $p_i^{\max}=0.25$; hard constraints use $p_i^{\max}=1$.}
    \label{tab:constraint_lims}
    \setlength{\tabcolsep}{4pt}
    \begin{tabular}{@{}llr@{}}
        \toprule
        \textbf{Constraint} & \textbf{Symbol} & \textbf{Limit} \\
        \midrule
        \multicolumn{3}{@{}l}{\textit{Soft}} \\
        Hip and thigh joint torque & $c_{\tau,hip}$   & 23.5 N$\cdot$m \\
        Knee joint torque          & $c_{\tau,knee}$  & 35.0 N$\cdot$m \\
        Leg joint velocity         & $c_{\dot{q}}$    & 30.0 rad/s \\
        Base linear velocity       & $c_v$            & 0.25 m/s \\
        Leg joint deviation        & $c_{dev}$        & 0.4 rad \\
        Foot stumble force ratio   & $K_f$            & 4 \\
        Foot slippage velocity     & $v_{\text{slip}}$& 0.2 m/s \\
        \addlinespace
        \multicolumn{3}{@{}l}{\textit{Hard}} \\
        Non-foot contact force & $c_{\text{contact}}$ & 1.0 N \\
        Body roll / pitch      & $c_{\text{fall}}$    & 90 deg \\
        Foot impact force      & $c_{f,\max}$         & 150.0 N \\
        \bottomrule
    \end{tabular}
\end{table}

\section{Training Configuration}
\setcounter{table}{0}
\renewcommand{\thetable}{D.\arabic{table}}

\Cref{tab:training} summarises the policy architecture, the control and simulation
rates, the action scaling of \cref{eq:action-scaling}, and the
hyperparameters of the constrained PPO implementation.

\begin{table}[H]
    \centering
    \caption{Training configuration}
    \label{tab:training}
    \setlength{\tabcolsep}{4pt}
    \begin{tabular}{@{}lr@{}}
        \toprule
        \textbf{Parameter} & \textbf{Value} \\
        \midrule
        \multicolumn{2}{@{}l}{\textit{Policy and control}} \\
        Hidden layers            & 512, 256, 128 \\
        Activation               & ELU \\
        Policy rate              & 50 Hz \\
        PD controller rate       & 200 Hz \\
        Episode duration         & 10 s \\
        Leg action scale         & 0.25 \\
        Arm action scale         & 0.50 \\
        \addlinespace
        \multicolumn{2}{@{}l}{\textit{Constrained PPO}} \\
        Parallel environments    & 4096 \\
        Batch size               & 98\,304 (4096 $\times$ 24) \\
        Mini-batch size          & 24\,576 (4096 $\times$ 6) \\
        Epochs                   & 5 \\
        Clip range               & 0.2 \\
        Entropy coefficient      & 0.005 \\
        Discount factor $\gamma$ & 0.99 \\
        GAE discount factor      & 0.95 \\
        Target KL divergence     & 0.01 \\
        Learning rate            & $3\times10^{-4}$ (adaptive) \\
        Training iterations      & 10\,000 \\
        \bottomrule
    \end{tabular}
\end{table}